\section{Introduction}
Test-time scaling (TTS) changes how a language model uses inference computation: it can sample answers, revise a draft, search intermediate states, or aggregate proposals. Transferring these mechanisms requires specifying what constitutes an answer, how states are evaluated, and which output is scored. These decisions determine whether a benchmark measures candidate generation, output selection, or a complete inference procedure.

Molecular instruction following provides a concrete setting for these choices. A model edits a molecule according to a request involving several properties. Many structures may satisfy the request, and improving one property need not improve the others. \gellmo{} supplies instruction-tuned checkpoints and tasks, with a published evaluation of twenty candidates per input \citep{gellmo2025}.

Declaring success whenever any candidate succeeds measures \emph{coverage}; testing the returned molecule measures \emph{terminal success}. High coverage need not imply successful selection. A selector with objective access can retrieve a successful candidate that likelihood ranking misses. A shared output count therefore requires explicit selection rules and disclosure of inference-time tools.

We compare nine strategies on two frozen \gellmo{} checkpoints, ten tasks, and all 3,905 task--input pairs per configuration. Every strategy returns one molecule ($K=1$), retaining its declared internal generation and selection rules. The comparison concerns complete procedures, with method-specific budgets and property-tool use.

Our contributions are a common terminal-output protocol that retains failures, a complete paired comparison documenting molecular adaptations of published mechanisms, and three complementary analyses: task success (Table~\ref{tab:effectiveness}), generation workload (Table~\ref{tab:workload}), and final-molecule quality (Table~\ref{tab:quality}). Cached-output diagnostics identify a substantial selection effect, while magnitude-sensitive evaluation changes the comparison between Sampling and Graph of Thoughts.
